\begin{corollary}[强局域 geometric uplift 下八重 parity 超选的 \texorpdfstring{$(4+4)$}{(4+4)} 坍缩]\label{cor:conclusion-window6-geometric-uplift-four-four-collapse}
沿用定理 \ref{thm:conclusion-window6-boundary-parity-three-layer-blind-filtration} 的记号，设
\[
P_6\cong (\ZZ/2\ZZ)^3,
\qquad
P_6^{\mathrm{geo}}=\langle(1,1,1)\rangle\cong \ZZ/2\ZZ.
\]
则对任意 fold-compatible 复表示
\[
M=\bigoplus_{\chi\in \widehat P_6} M_\chi
\]
的八重 parity 超选分解，限制到 \(P_6^{\mathrm{geo}}\) 后都精确坍缩为两类，每类恰含四个角色扇区：
\[
M=M_{\mathrm{geo},+}\oplus M_{\mathrm{geo},-},
\qquad
M_{\mathrm{geo},\pm}
=
\bigoplus_{\chi|_{P_6^{\mathrm{geo}}}=\pm 1} M_\chi.
\]
等价地，限制映射
\[
\widehat P_6\longrightarrow \widehat P_6^{\mathrm{geo}} \cong \{\pm 1\}
\]
的两个纤维都恰有四个元素。
\leanverified{paper\_conclusion\_geometric\_uplift\_four\_four\_collapse}
\end{corollary}
\begin{proof}
由定理 \ref{thm:conclusion-window6-boundary-parity-three-layer-blind-filtration}，
\[
P_6\cong (\ZZ/2\ZZ)^3,
\qquad
P_6^{\mathrm{geo}}=\langle(1,1,1)\rangle.
\]
写
\[
\widehat P_6\cong (\ZZ/2\ZZ)^3\ni a=(a_1,a_2,a_3).
\]
则角色 \(\chi_a\) 在几何对角元上的取值为
\[
\chi_a(1,1,1)=(-1)^{a_1+a_2+a_3}.
\]
因此恰有四个 \(a\) 满足 \(a_1+a_2+a_3\equiv 0\pmod 2\)，另四个满足奇 parity。于是限制映射的两个纤维都含四个元素，从而八重超选精确坍缩为 \((4+4)\) 二分。证毕。
\end{proof}

\begin{theorem}[window-\texorpdfstring{$6$}{6} boundary parity 的三重外置性]\label{thm:conclusion-window6-boundary-parity-threefold-externality}
window-\(6\) 的 canonical boundary parity 同时满足以下三条彼此独立的性质：
\begin{enumerate}
  \item \textbf{表示可见：}它在表示论层面给出严格的八重超选分解；
  \item \textbf{有理连续不可见：}它在 \(\QQ\)-系数的 de Rham/Sullivan/有理特征类机制下完全不可见；
  \item \textbf{强局域几何不完备：}在强局域 geometric uplift 下，只剩对角子群
  \[
  P_6^{\mathrm{geo}}=\langle(1,1,1)\rangle\cong \ZZ/2\ZZ
  \]
  的粗粒影像，余下两位 parity 不能由真几何局域机制承载。
\end{enumerate}
因此，boundary parity 不是某种可以由单一连通连续证书、单一有理包络或单一局域几何模型完全吸收的附属数据；它本质上是一份三重外置的 \(2\)-初离散记忆。
\end{theorem}
\begin{proof}
第 \(1\) 条由定理 \ref{thm:conclusion-window6-boundary-superselection-c3-orbit-stratification} 给出：boundary 中心
\[
P_6\cong (\ZZ/2\ZZ)^3
\]
强制出现八个角色扇区，并在 \(C_3\)-pinning 下形成严格超选分解。

第 \(2\) 条由定理 \ref{thm:conclusion-window6-boundary-parity-directsummand-rational-blindness} 直接给出：对任意连通 Lie 包络，boundary parity 在正次数有理上同调上的拉回恒为零，故它对一切 \(\QQ\)-系数连续特征类完全不可见。

第 \(3\) 条由定理 \ref{thm:conclusion-window6-boundary-parity-three-layer-blind-filtration} 与推论 \ref{cor:conclusion-window6-geometric-diagonal-fourfold-residual-cover} 给出：强局域 geometric uplift 唯一可实现的方向只是
\[
P_6^{\mathrm{geo}}=\langle(1,1,1)\rangle,
\]
从而八重超选只能坍缩为两类四重几何残差，而不能在真几何层面恢复全部三位 parity。证毕。
\end{proof}

\begin{theorem}[抽象中心恢复 boundary cube 的十七比特外置信息壁垒]\label{thm:conclusion-window6-boundary-center-seventeen-bit-barrier}
设
\[
Z_\partial\cong (\ZZ/2\ZZ)^3\hookrightarrow Z(G)\cong (\ZZ/2\ZZ)^8
\]
为 window-\(6\) 的 canonical boundary 中心三维子空间。则
\[
\bigl|\operatorname{Orb}_{\operatorname{Aut}(G)}(Z_\partial)\bigr|
=
\binom{8}{3}_2
=
97155.
\]
因此，仅从抽象群 \(G\) 出发，若要精确恢复哪一个三维中心子空间对应 boundary parity，至少需要
\[
\left\lceil \log_2 97155\right\rceil = 17
\]
个二值选择；其 Shannon 信息量恰为
\[
\log_2 97155 \approx 16.56800062435218.
\]

另一方面，visible 根云的矩张量 pinning 只剩一个 Klein 四量级的残余对称：
\[
\operatorname{Stab}_{O(3)}(\mathfrak M_6,\mathfrak m_6)\cong (\ZZ/2\ZZ)^2,
\]
故至多对应两位二值歧义。于是 hidden 抽象中心恢复与 visible 根云 pinning 之间存在至少
\[
17-2=15
\]
比特的严格信息不对称。
\leanverified{paper\_conclusion\_boundary\_center\_seventeen\_bit\_barrier}
\end{theorem}
\begin{proof}
由命题 \ref{prop:conclusion-window6-ramanujan-defect-projection-geometric-not-abstract}，\(\Aut(G)\) 在
\[
Z(G)\cong \FF_2^8
\]
的全部三维子空间上作用传递；同一命题明确给出其轨道大小恰为
\[
\binom{8}{3}_2=97155.
\]
故若无额外几何 pinning，抽象群层面无法区分这 \(97155\) 个候选三维子空间。要逐一编码它们，至少需要
\[
\lceil \log_2 97155\rceil = 17
\]
个二值决策，而 Shannon 信息量即为 \(\log_2 97155\)。

另一方面，推论 \ref{cor:conclusion-window6-moment-tensor-axis-pinning} 已证明 visible 根云的联合稳定子满足
\[
\operatorname{Stab}_{O(3)}(\mathfrak M_6,\mathfrak m_6)\cong (\ZZ/2\ZZ)^2,
\]
故 visible 侧只剩两位二值歧义。二者相减便得到至少 \(15\) 比特的信息不对称。证毕。
\end{proof}

\begin{corollary}[连通动力学对称与几何 Lie 机制的交只剩环面相位影]\label{cor:conclusion-window6-torus-phase-only-common-shadow}
设 \(H\) 为作用于
\[
H_6=\ell^2(X_6,\pi_6)\cong \CC\langle X_6\rangle
\]
的连通紧 Lie 群，且其表示 \(\rho\) 与 pushforward 动力学算子 \(\mathsf T_6\) 对易。则经适当酉共轭后
\[
\rho(H)\subseteq U(1)^{21},
\]
从而 \(\operatorname{Lie}(\rho(H))\) 必为阿贝尔 Lie 代数。特别地，window-\(6\) 的几何 \(B_3/C_3\) Lie 机制不可能作为动力学中心化子的一部分出现；二者的共同连通影像至多是环面相位，而不可能保留任何非阿贝尔 Chevalley 结构。
\end{corollary}
\begin{proof}
由推论 \ref{cor:conclusion-window6-dynamical-centralizer-toralization}，任何与 \(\mathsf T_6\) 对易的紧 Lie 群都可酉共轭入
\[
U(1)^{21},
\]
故其 Lie 代数必阿贝尔。

另一方面，定理 \ref{thm:conclusion-window6-dual-lie-realizations-incompatible} 已证明：同载体上的几何 Lie 结构与动力学 Lie 包络 \(\mathfrak{sl}_{21}\) 彼此不相容，不能由同一组内导子统一实现。故二者的共同连通影像只能退化为环面相位，而不可能保留任何非阿贝尔 \(B_3/C_3\) 结构。证毕。
\end{proof}

\begin{conjecture}[window-\texorpdfstring{$6$}{6} 的 Walsh--Zeckendorf 对偶]\label{conj:conclusion-window6-walsh-zeckendorf-duality}
在刚性识别
\[
S_{6,2}\cong X_4\cong \{0,1,\dots,7\},
\qquad
\widehat Z_6^{\mathrm{bd}}\cong (\ZZ/2\ZZ)^3
\]
下，存在一个 fold-compatible 的有限 Fourier--Walsh 变换
\[
\mathcal W_6:\CC[S_{6,2}]\longrightarrow \CC[\widehat Z_6^{\mathrm{bd}}],
\]
满足：
\begin{enumerate}
  \item Zeckendorf 终端窗口 \(13+\{0,\dots,7\}\) 的坐标基被送到八个 boundary parity 角色基；
  \item 强局域几何可实现的对角子群 \(\langle(1,1,1)\rangle\) 在频域上恰对应推论 \ref{cor:conclusion-window6-geometric-uplift-four-four-collapse} 的 \((4+4)\) 粗粒二分；
  \item 最小退化壳中的五个 cyclic 极小点与三个 boundary 零权点，在 \(\mathcal W_6\) 下分别对应非平凡 Walsh 模与三条坐标角色。
\end{enumerate}
若此猜想成立，则 window-\(6\) 的最小退化壳、boundary parity 立方体与 Walsh 角色层并非三套平行字典，而是同一有限 Fourier 对偶的三种坐标化。
\end{conjecture}
\begin{proof}[证据]
定理 \ref{thm:conclusion-window6-minimal-shell-rigid-subcover-root-slice} 与定理 \ref{thm:conclusion-window6-residual-threeinterval-terminal-window} 已把
\[
S_{6,2}
\]
刚性识别为一个长度为 \(8\) 的终端 Zeckendorf 窗；推论 \ref{cor:conclusion-window6-geometric-uplift-four-four-collapse} 则表明八个 boundary parity 角色在强局域几何下精确坍缩为 \((4+4)\) 二分。另一方面，定理 \ref{thm:conclusion-window6-walsh-spin-duality} 已经把 window-\(6\) 的 boundary 角色层与 Walsh/Spin 词典联系起来。

再结合定理 \ref{thm:conclusion-window6-minimal-shell-five-three-cleavage} 的 \((5+3)\) 最小壳裂分，可以看出：终端 Zeckendorf 窗、八点 parity 角色立方体以及 Walsh 模态层之间，已经存在一组与有限 Fourier 对偶高度一致的结构线索。现阶段文稿尚未给出 \(\mathcal W_6\) 的显式构造，故保留为猜想。证毕。
\end{proof}
